\subsection{Connection through a private store}
\label{sec:private-connection}

The public boundary of \cref{sec:composition} is useful when the next
component should see the whole execution. A different interface retains
the store and its audiences, including private cells. It gives a
connection rule for the fixed-schedule disclosure language without
requiring publication at the boundary.

\paragraph{Open fragments.}
Replace $\mathsf{done}$ by a hole in a finite disclosure protocol.
A fragment judgment
\[
 P:(\Xi,\Omega)\longrightarrow(\Xi',\Omega')
\]
processes its commands using exactly the phase rules of
\cref{sec:disclosure}, but permits pending cells at both ends. The
output context retains all input declarations in their original order,
with the same types and owners; only audiences can grow. New binds
have fresh names. The hole makes no choice and publishes no record.
The connection rule is
\[
 \frac{P:(\Xi,\Omega)\longrightarrow(\Xi',\Omega')\qquad
       Q:(\Xi',\Omega')\longrightarrow(\Xi'',\Omega'')}
      {P;Q:(\Xi,\Omega)\longrightarrow(\Xi'',\Omega'')},
\]
where $P;Q$ fills the hole of $P$ with $Q$, after fresh bound names
have been chosen. Interfaces match as ordered dependent contexts,
up to consistent renaming and data conversion. No value is copied to
a new witness at this boundary. A complete protocol additionally has
empty final pending set.

\paragraph{All possible boundary stores.}
Fix a set-theoretic model of the data calculus and a family $W$ of
well-typed input stores for $\Xi$. Let $W_m(P;W)$ be the stores
reachable after $m$ commands when every bind may choose any value of
its instantiated type. This definition is independent of a chosen
profile. Write $W_P$ for the final family. Publications change only
the phase metadata. Immutable cells and the fixed schedule mean that
a store at phase $m$, together with that schedule, determines its
entire preceding history from the input interface: earlier stores
are restrictions and earlier audiences are known. Each local view
likewise determines all the role's preceding local views.

At a bind occurrence $m$ owned by $i$, put
\[
 \mathcal V_m=\{V_i^m(\rho):\rho\in W_m(P;W)\}.
\]
Here $m$ counts commands preceding the bind. Its answer type is a
family $A_m(v)$ on $\mathcal V_m$, since its dependencies are in
the actor's local context. Define the full strategy space
\begin{equation}\label{dep:eq:fragment-strategies}
 \mathsf S_i(P;W)=
 \prod_{\substack{m\text{ a bind in }P\\\text{owned by }i}}
       \prod_{v\in\mathcal V_m} A_m(v).
\end{equation}
Every legal view receives a decision, including views inconsistent with
the player's own earlier decisions. An empty product is a singleton;
no inhabitance claim is made if a required answer type is empty.
A data implementation induces such a strategy by evaluation, whereas
the set-theoretic space also admits nondefinable functions, as in
\cref{sec:policies}. Smaller implementation supports certify dependency
claims; they do not silently restrict the full space of deviations in
\eqref{dep:eq:fragment-strategies}.

\begin{theorem}[Private-store connection]\label{dep:thm:connection}
For well-formed fragments $P,Q$ with matching boundary and input family
$W$, connection has the following properties.
\begin{enumerate}
\item Every execution preserves dependent typing and each stored witness.
Its boundary store belongs to $W_P$, and
\[
 \mathsf{Run}_{P;Q}(\sigma,\tau;\rho)
 =\mathsf{Run}_Q(\tau;\mathsf{Run}_P(\sigma;\rho)).
\]
The stores retain all intermediate bound values.
\item There is a canonical bijection
\[
 \mathsf S_i(P;Q;W)\cong
 \mathsf S_i(P;W)\times\mathsf S_i(Q;W_P).
\]
It preserves every decision at every legal view, and carries whole-player
unilateral deviations to pairs of component deviations.
\item Connection is associative, with the empty fragment as unit,
up to consistent bound-name renaming and the displayed strategy
bijections. Rebracketing preserves histories, dependent answer families,
observation partitions, and execution.
\item Sealing certificates compose when the marked dependency set at
the end of the first fragment is passed unchanged to the second.
The two-execution hiding theorem therefore holds across this private
boundary for implementations satisfying those certificates.
\end{enumerate}
\end{theorem}
\begin{proof}
The second fragment starts with exactly the context, pending set, and
store produced by the first. Phase preservation applies at every
command. Fresh binding and immutable storage preserve all old values;
publication resolves the same named witness. This proves the first
claim and the execution equation by induction on $P$.

All stores in $W_P$ are endpoints of legal $P$ histories. Appending any
legal $Q$ history to one of those histories gives a legal composite
history; conversely every composite history splits there. Therefore
the stores at each $Q$ phase are exactly those obtained by starting
$Q$ from $W_P$. The boundary has no extra observation event. Its
local context already contains every value that the role previously
saw, so the suffix's view agrees with the composite view, with only
the fixed phase offset changed. The factors of
\eqref{dep:eq:fragment-strategies} consequently split into the binds
in $P$ and the binds in $Q$. Splitting and concatenating these decision
families are inverse. This also proves the claim about deviations;
the suffix is indexed by all legal boundary stores, not merely the
one reached by the chosen prefix strategy.

For three fragments, both bracketings are the same command list and
have the same families of reachable stores at each phase. In
particular $(W_P)_Q=W_{P;Q}$. The strategy identifications are just
reassociation of three products. The empty list has identity execution
and singleton strategy spaces. Because audiences and answer families
at every phase are unchanged, the history and observation claims
follow as well.

Finally, the sealing scan is a deterministic update on the marked
set. Running it over $P;Q$ is running it first over $P$ and then
over $Q$ with the resulting marks. At the boundary the two runs
agree on all unmarked cells, and all marked cells remain invisible
to the observer. These are exactly the induction invariants of
\cref{thm:hiding}; no reset of the marked set or of the store occurs.
Its proof therefore continues through the boundary.
\end{proof}

For example, a first fragment may privately bind $n:\Nat_{>0}$ and
$v:\mathsf{Vec}(\Bool,n)$, and a second may let Bob bind an independent
Boolean before Alice publishes $n$ and then $v$. Both pending cells
cross the boundary. Bob's strategy before publication has no argument
for either value; after publication the vector still has the original
index. Publishing $v$ first remains ill typed. The interface retains
both the dependent witness and the information restriction.

This settles connection for this finite, fixed-schedule phase
language. It does not give a cut rule for asynchronous channels,
data-dependent control, or observable failure. It composes full games and implementations;
no equilibrium set is selected at an intermediate boundary. Thus it
supplies no marginal-equilibrium composition law of the refuted kind
in \cref{se:sec:open-game-obstruction}.
